\color{unimain}{\subsection{Орган выявления БНТ}}
\color{black}

\begin{enumerate}


\item
Для обеспечения несрабатывания МТЗ, ТЗНП и ТЗОП при броске тока намагничивания используется блокировка по содержанию второй гармоники в токе ВВ. Блокировка основана на контроле отношения действующего значения второй гармоники к первой тока фаз. Имеется возможность регулировать порог срабатывания блокировки, с помощью уставки \textbf{«Ih2/Ih1 БНТ»}. Блокировка по второй гармонике выводится из действия при величине тока ниже 0,04Iном, при срабатывании измерительного органа \textbf{«Iср БНТ»}, а так же при выводе программной накладкой \textbf{«БНТ ВВ»}.

\item 
Уставки ОВ БНТ приведены в таблице \ref{app_set:ovbnt}. Алгоритм работы ОВ БНТ приведен на рисунке \ref{pict:BNT}.

\begin{figure}[H]
\centering
\includegraphics[width=1\textwidth,height=1\textheight,keepaspectratio]{\fbpath/09.01 БСК/271.0391 ОВ БНТ ВВ/img/BNT.pdf}
\captionof{figure}{Функциональная схема логики ОВ БНТ}\label{pict:BNT}
\end{figure}


\end{enumerate} 